%---------------------------------------------------------
% PART 2 — MATHEMATICAL FRAMEWORK AND DERIVATIONS
%---------------------------------------------------------

\section{Mathematical Framework of UST}

Unified Source Theory (UST) is built on the hypothesis that a single
dimensionless constant $N_{s,q}$ governs stabilisation phenomena across
quantum information theory, gravitational geometry, and cosmology. In this
section we present the mathematical foundations of UST v3, beginning with the
geometric derivation of $N_{s,q}$, followed by its quantum-informational
interpretation and its role in the Omnium field formalism.

%---------------------------------------------------------
\subsection{Geometric Derivation of the Universal Constant $N_{s,q}$}

The geometric origin of $N_{s,q}$ arises from the extremization of the
Einstein tensor component $G_{tt}$ for a static, spherically symmetric metric
of the form
\[
ds^2 = -A(r)\,dt^2 + A(r)^{-1}dr^2 + r^2 d\Omega^2.
\]
The Einstein tensor for this metric yields
\[
G_{tt} = A(1-A) + rA',
\]
where $A' = dA/dr$. At the Omnium horizon $r=r_s$, UST imposes the condition
that the geometric contribution to the energy density is maximized. Setting
$A'=0$ at the extremum gives
\[
G_{tt}(r_s) = A(1-A).
\]
The extremum of the function $f(A)=A(1-A)$ occurs at
\[
\frac{df}{dA} = 1 - 2A = 0 \quad \Rightarrow \quad A = \frac{1}{2}.
\]
UST identifies the horizon ratio as
\[
A = (1-N_{s,q})(2-N_{s,q}),
\]
which yields the quadratic equation
\[
(1-N_{s,q})(2-N_{s,q}) = \frac{1}{2}.
\]
Expanding and solving,
\[
2 - 3N_{s,q} + N_{s,q}^2 = \frac{1}{2},
\]
\[
N_{s,q}^2 - 3N_{s,q} + \frac{3}{2} = 0.
\]
The physically relevant root is
\[
N_{s,q} = \frac{3 - \sqrt{3}}{2} = 0.63354460.
\]
This value forms the geometric backbone of UST.

%---------------------------------------------------------
\subsection{Quantum-Informational Interpretation}

In UST, the observable universe (Channel $Q$) and a frozen, non-interacting
sector (Channel $C$) form a composite system whose stabilisation is governed
by $N_{s,q}$. The decoherence-free subspace (DFS) formalism yields the
stabilisation condition
\[
K \cdot dt = \pi N_{s,q},
\]
where $K$ is the stabilisation rate and $dt$ is the characteristic time
interval of the DFS. This relation is derived from the fidelity bound
\[
F(t) = \cos^2(Kt),
\]
whose minimum stabilisation requirement is
\[
Kt = \frac{\pi}{2}N_{s,q}.
\]
Thus the universal stabilisation constant emerges naturally from the DFS
structure.

%---------------------------------------------------------
\subsection{The Omnium Field and Two-Channel Structure}

UST postulates that the total energy density of the universe is given by the
Omnium field
\[
\rho_{\rm Om} = N_{s,q}\rho_Q + (1-N_{s,q})\rho_C,
\]
where $\rho_Q$ is the energy density of the observable sector and $\rho_C$ is
the energy density of the frozen sector. The ratio
\[
\frac{N_{s,q}}{1-N_{s,q}} = 1.7288
\]
matches the SU(3) color ratio, providing a group-theoretic interpretation of
the two-channel structure.

Channel $C$ is interpreted as the source of dark matter and dark energy:
\[
\rho_{\rm DM} \propto (1-N_{s,q}), \qquad
\rho_{\rm DE} \propto N_{s,q}.
\]

%---------------------------------------------------------
\subsection{Stabilisation Constant and the $\pi N_{s,q}$ Relation}

The stabilisation constant arises from the requirement that the Omnium field
remains dynamically stable under perturbations. The effective Hamiltonian for
the two-channel system is
\[
H_{\rm eff} = N_{s,q}H_Q + (1-N_{s,q})H_C.
\]
The stabilisation condition is obtained by requiring that the phase difference
between the two channels remains bounded:
\[
\Delta\phi = \int (E_Q - E_C)\,dt = \pi N_{s,q}.
\]
Thus the universal stabilisation relation
\[
K\,dt = \pi N_{s,q}
\]
is a direct consequence of the two-channel structure.

%---------------------------------------------------------
\subsection{Summary of Mathematical Foundations}

The mathematical framework of UST v3 rests on the following pillars:
\begin{itemize}
\item Geometric extremization of the Einstein tensor yields $N_{s,q}$.
\item Quantum-informational stabilisation leads to $K\,dt = \pi N_{s,q}$.
\item The Omnium field formalism unifies the observable and frozen sectors.
\item The ratio $N_{s,q}/(1-N_{s,q})$ matches SU(3) structure.
\item Channel $C$ provides a natural explanation for dark matter and dark energy.
\end{itemize}

These foundations support the observational and theoretical developments
presented in subsequent sections.